% Source: https://artofproblemsolving.com/community/c3h1726318_prove_consecutive_integers
% Stable identity: aops-topic-1726318-q01
\textbf{Sums of consecutive integers.} A sum of $k\ge1$ consecutive positive integers starting at $a\ge1$ is
\[
a+(a+1)+\cdots+(a+k-1)=\dfrac{k(2a+k-1)}{2}.
\]
So such a sum equals $N$ exactly when $2N=k\,m$ with $m=2a+k-1$. Here $m-k=2a-1$ is odd and positive, so $k$ and $m$ have \emph{opposite parity} and $m>k$. Conversely, if $2N=km$ with $k<m$ of opposite parity, then $a=\dfrac{m-k+1}{2}$ is a positive integer, and the $k$ integers from $a$ to $a+k-1$ add up to $N$. So the representations of $N$ as a sum of $k\ge1$ consecutive positive integers correspond exactly to the factorisations $2N=km$ with $k<m$ and $k,m$ of opposite parity.

\textbf{Counting the factorisations.} In such a factorisation exactly one of $k,m$ is odd, and it divides $N$ (since the odd factor $d$ of $2N=dd'$ is coprime to $2$). Conversely, every odd divisor $d$ of $N$ gives the factorisation $2N=d\cdot\dfrac{2N}{d}$ into an odd and an even number, which are different; putting the smaller one first gives a pair $k<m$. Different odd divisors give different pairs (the odd member of the pair is $d$). Hence the number of representations with $k\ge1$ terms equals the number $n$ of odd divisors of $N$.

\textbf{Removing the trivial sum.} Exactly one of these representations has $k=1$ term, namely $N$ itself (it comes from $d=1$ if $N$ is even, and from $d=N$ if $N$ is odd). So $N$ has exactly $n-1$ representations as a sum of \emph{two or more} consecutive positive integers. $\blacksquare$

For example, $15$ has odd divisors $1,3,5,15$, and indeed $15=7+8=4+5+6=1+2+3+4+5$ in $3$ ways; a power of $2$ has $n=1$ and no such representation.
